\documentclass[12pt]{article}
\usepackage{Style-Thesis-Report}

% You can define your own commands like this
\newcommand{\um}{$\upmu$m} % easily creates the label for ``micrometers''





\begin{document}
	
\begin{titlepage}
	\begin{center}
		
		\vspace*{1cm}
		
		\Huge
		\textbf{pooop Determining the Crystaline Structure of Metals through X-Ray Diffraction}
		\large
		\vspace{1.5cm}
		
        Shukri Alex, Weltz Michael
		
		McGill University Department of Physics
		
		Supervisor: Brandon
		
		\today
		
	\end{center}
	
\end{titlepage}

\begin{abstract}

	% put some abstract stuff in the into. From journal-template, the theory is more relevant in intro



This experiment set out to to measure X-ray diffraction patterns of various solid 
samples in order to describe their crystal structure at the atomic/molecular level. An 
already calibrated Siemen !!!!!MODEL NAME!!!!! X-Ray diffractometer was used to send X-rays 
into Copper (Cu), Nickel (Ni), CuNi alloy, Titanium (Sn), Lead (Pb), PbSn alloy and two
mystery samples named A and B. 



Any crystalline solid consists of a regular 3-dimensional packing of cells, the shape of 
which corresponds to one of fourteen Bravais lattices. These 3D patterns create multiple 
parallel planar layers within the crystal of various orientations.
X-Rays travelling through these planes will diffract off the atoms and constructively 
interfere at the detector, creating a peak in photon count in the observed diffraction pattern. 
From that peak, Bragg's law $n\lambda = 2d \sin \theta$ (where $n=1$ for a first-order plane) can be used to 
deduce the inter-planar distance $d$ of the diffraction plane of the crystal. Using this 
inter-planar distance and the corresponding Miller indices, the Bravais lattice can be determined
and the lattice parameter(s) calculated. The Cu and Ni samples were found to have a Bravais 
lattice of [ ... ], with parameters [ ... ] respectively. The alloy sample of these two 
was found to have [ ... ]. The Sn and Pb samples were found to have a lattice of [...] 
and parameters [...] respectively, and their alloy [...], [...]. Finally, the mystery 
samples A and B were determined to be [...] and [...], with lattices [...] and parameters
[...], respectively. 


\end{abstract}
\tableofcontents{}

\thispagestyle{empty}
\pagebreak
\setcounter{page}{1}


\section{Introduction}\label{sec:int}

The discovery ofX-rays led to W.H. and W.L. Bragg researching the inner 
workings of X-ray diffraction. It was found a solid can serve as a diffraction
grating for x-rays of similar lengths to the inter-atomic distances of the solid.
Similarly to how light diffracts, one can obtain a relation between the inter-atomic
distance $d$, the incident angle $\theta$, and the wavelength $\lambda$: 
\begin{equation}\label{eq:bragg}
	n\lambda=2d\sin(\theta)
\end{equation}

with $n=1$ for our purposes. 



\section{Methods}\label{sec:met}

A pre-calibrated Siemens diffractometer was used to produce X-Rays, send them at the samples, and 
detect the diffractions.
The diffractometer consisted primarily of an X-ray tube, a sample mounting plate, and a 
photon detector. The X-ray tube used Copper $K_{\alpha}$ radiation and was cooled using water. Two sets of filters of 
0.6\,mm were present on the X-ray tube and the detector, to prevent stray rays from leaving
the tube and entering the detector. During data acquisition, the diffractometer was set to 
40\,kV and 40\,mA. The mounting plate kept the sample in place with a spring-loaded mechanism. 
The samples were centered by aligning the middle of the sample with a red line marked on the 
moutning plate. The diffractometer changes the $\theta$ by rotating the mounting plate by some $\Delta\theta$
and moving the detector by $2\Delta\theta$.
A computer was connected to the diffractometer, with which
a .dql measuring parameter file was given to the diffractometer and data was obtained from the 
dectector as .RAW files. The measuring parameters were an initial angle of $2\theta_0=5.00^{\circ}$, 
angle increments of $\Delta\theta=0.05^{\circ}$, with 2.00\,s at each step, with a final 
angle ofe $2\theta_f=110.00^{\circ}$. The total measuring time is roughly 1\,h0\,5m using the aformentionned
parameters.
Data acquisition was simple: Mount the desired sample onto the mounting plate, centering it
with the red line. Then, on the computer, open the software program [!!!NAME!!!], select the desired 
parameter file, name the trial, and execute the job. After the job finishes, convert the .RAW file to .UXD.
Finally, the .UXD files were transfered from the lab computer onto the experimentalists's for 
the data analysis.
Three trials were done for Ni and Cu samples, though the data for the first Cu trial did not save correctly.
Two trials were done for Pb, Sn, A and B. For each of the alloys (including different percentage of mixing) 
one trial was done. It was determined that the uncertainty due to random errors would be 
less that due to noise (could expand on this).




\section{Results}\label{sec:res}

From the .UXD files of a sample, the plot of light intensity against $\theta$ can be done.
Figure \ref{} show this plot for the Copper and Nickel samples. Using prominence 50, 
local maxima of light intensities can be located and indexed to find angles where diffraction in planes
occur. Every peak in the plot is then fitted to a gaussian curve, centered at a local maxima,
and every curve is summed to obtain a plot of summed gaussians of
intensity against $2\theta$. Figure \ref{} shows the summed gaussians for Copper and 
Nickel. The mean angle for every peak as well as its uncertainty can then be obtained from
the corresponding gaussian. Table










\section{Discussion}\label{sec:dis}

\section{Conclusion}\label{sec:con}




% Bibliography:
\bibliographystyle{IEEEtran}
\bibliography{MyBibliography}

\end{document}